\documentclass[11pt,a4paper]{article}
\usepackage[utf8]{inputenc}
\usepackage[T1]{fontenc}
\usepackage{lmodern}
\usepackage{textcomp}
\usepackage{amsmath,amssymb}
\usepackage{graphicx}
\usepackage{booktabs}
\usepackage{longtable}
\usepackage{array}
\usepackage{listings}
\usepackage{xcolor}
\usepackage[margin=1in]{geometry}
\usepackage{float}
\usepackage[hidelinks,breaklinks]{hyperref}
\usepackage{microtype}
\setlength{\parskip}{4pt}
\lstset{basicstyle=\ttfamily\scriptsize,breaklines=true,breakatwhitespace=false,
  postbreak=\mbox{\textcolor{gray}{$\hookrightarrow$}\space},frame=single,
  framerule=0.3pt,rulecolor=\color{gray!60},columns=fullflexible,keepspaces=true,
  showstringspaces=false,numbers=none,xleftmargin=4pt,xrightmargin=4pt,
  literate={\_}{{\_}}1}
\graphicspath{{figures/}}
\providecommand{\tightlist}{}
\title{Supplementary material: The Hubble Tension: A Conceptual Summary Prepared Without Computation or Literature Search}
\author{HEP-Genesis-Agent session record \\ \small Project: ramp-features}
\date{2026-10-16}
\begin{document}
\maketitle
\noindent\textit{\small Supplementary material, generated mechanically from the session record: it lists exactly what ran, in order, without interpretation.}
\vspace{1em}



\appendix
\section{Implementation details}
\label{app:impl}
This appendix is generated mechanically from the session record; it lists exactly what ran, in order, without interpretation.
\subsection{Session}
\begin{itemize}
  \item Source: cli client; conversation \texttt{1dffc244}
  \item Started: 2026-10-15 22:42:59; report generated: 2026-10-16 17:02:56
  \item Project: ramp-features (\texttt{/\allowbreak{}Users/\allowbreak{}nramachandra/\allowbreak{}hep-genesis-testing/\allowbreak{}projects/\allowbreak{}ramp-features})
  \item Compute execution site(s) selected in the session: local
  \item 2 user turn(s), 1 tool call(s) (0 returned errors), 0 figure(s), 0 data file(s)
\end{itemize}
\subsection{Models}
\begin{longtable}{@{}p{0.16\textwidth}p{0.34\textwidth}p{0.44\textwidth}@{}}
\toprule Stage & Model & Endpoint \\ \midrule \endhead
quick & \texttt{claude-opus-5-5} & \texttt{https:\allowbreak{}/\allowbreak{}/\allowbreak{}mag.genesis.americansciencecloud.org/\allowbreak{}v1} \\
report & \texttt{claude-opus-5-5} & \texttt{https:\allowbreak{}/\allowbreak{}/\allowbreak{}mag.genesis.americansciencecloud.org/\allowbreak{}v1} \\
coder & \texttt{claude-opus-5-5} & \texttt{https:\allowbreak{}/\allowbreak{}/\allowbreak{}mag.genesis.americansciencecloud.org/\allowbreak{}v1} \\
\bottomrule\end{longtable}
\subsection{Tool-call log}
Chronological; arguments and results are truncated for space (the full record is in the project's \texttt{provenance/} folder).
\begin{longtable}{@{}p{0.04\textwidth}p{0.22\textwidth}p{0.36\textwidth}p{0.32\textwidth}@{}}
\toprule \# & Tool & Arguments & Outcome \\ \midrule \endhead
1 & \texttt{genesis}\newline\texttt{write\_\allowbreak{}report} & \texttt{\{''focus'':\allowbreak{} ''Conceptual explanation of the Hubble tension:\allowbreak{} early-universe (Planck CMB,\allowbreak{} arXiv:\allowbreak{}1807.06209,\allowbreak{} H0=\allowbreak{}67.4+/\allowbreak{}-0.5)\allowbreak{} vs late-universe distance-ladder (SH0ES,\allowbreak{} arXiv:\allowbreak{}2112.04510,\allowbreak{} H0=\allowbreak{}73.04+/\allowbreak{}-1.04)\allowbreak{} measurements,\allowbreak{} \textasciitilde{}\allowbreak{}5 sigma dis\ldots{}\allowbreak{}} & \texttt{Report generated:\allowbreak{} /\allowbreak{}Users/\allowbreak{}nramachandra/\allowbreak{}hep-genesis-testing/\allowbreak{}projects/\allowbreak{}ramp-features/\allowbreak{}reports/\allowbreak{}261015\_\allowbreak{}224313/\allowbreak{}report.pdf Title:\allowbreak{} The Hubble Tension:\allowbreak{} A Conceptual Summary of Early- and Late-Universe Determinat\ldots{}\allowbreak{}} \\
\bottomrule\end{longtable}
\end{document}
